\section{METHODOLOGY}\label{ch:method}

The tail risk of a portfolio depends on how volatile each asset is and on how the assets move together, and a market regime can change either. The models in this chapter share one regime signal, one portfolio and one estimation scheme, and differ in whether volatility, correlation or both depend on the regime, and in whether volatility also follows GARCH dynamics. Pairing them so that each comparison changes as little as possible (Section~\ref{sec:ladder-design}) is what lets the thesis ask about the correlation channel specifically. Chapter~\ref{ch:eval} defines the loss, the test and the decision rule, and Chapter~\ref{ch:code} gives the implementation settings.

\subsection{The forecasting task and the estimation scheme}\label{sec:walkforward}

For each forecast day $t$ I forecast the one-day-ahead value-at-risk $\VaR_t$ and expected shortfall $\ES_t$ of the portfolio return $y_t$ at tail level $\alpha$, using returns through day $t-1$ only. $\VaR_t$ is the $\alpha$-quantile of the forecast distribution of $y_t$, and $\ES_t$ is the mean of $y_t$ below that quantile. Both are negative numbers. The primary tail level is $\alpha=5\%$ and the secondary one is $2.5\%$.

Forecasts come from a walk-forward scheme. The forecast days are split into blocks of 21 consecutive days. For the block starting at day $t_0$, every model is estimated on the rolling window of the 1,260 returns before $t_0$ (about five years) and used, without re-estimation, for the 21 days of the block. The regime model is refitted at every block. The GARCH margins of Section~\ref{sec:garch} are refitted every third block, that is every 63 days, counting from the first block of the run; at the other blocks their parameters are kept and only the window is re-filtered. The first block of a run starts the hidden-state chain cold (Section~\ref{sec:hmm}). A model is therefore a forecasting procedure, consisting of the specification, the window length and the refit schedule, and this is the object that Chapter~\ref{ch:eval} compares.

\subsection{The regime signal}\label{sec:hmm}

The regime signal is a two-state Gaussian hidden Markov model (HMM) fitted to the SPY daily simple return only, so that the same signal feeds every model. Let $S_t\in\{1,2\}$ be the hidden state, with state 1 the calm state (lower return variance) and state 2 the stress state. Given $S_t=k$, the SPY return is normal with mean $\mu_k$ and variance $\sigma_k^2$, and $S_t$ is a Markov chain with transition matrix $P$. I estimate $(\mu_k,\sigma_k^2,P)$ on the window by the expectation--maximisation algorithm. The first fit of a run starts from fixed initial values, and every later fit warm-starts from the previous block's estimates. States are ordered by variance after each fit, so the stress state keeps its label.

Two sets of probabilities are used, and they must not be confused. Inside the window, the smoothed probabilities $\gamma_{s,k}=\Pr(S_s=k\mid\text{window})$ for each window day $s$ weight the regime moments below; they use data up to the end of the window and nothing after it. For the forecast, the weights of the mixture are predictive probabilities
\begin{equation}
p_{t,k}=\Pr\bigl(S_t=k\mid \text{SPY returns through } t-1\bigr),
\end{equation}
computed by propagating the filtered probability at the window end through $P$ and, within the block, filtering with the realised SPY returns up to day $t-1$. Day $t$'s own return is never used.

\subsection{Static regime-covariance models: A, B and D}\label{sec:static}

For weights $q_s\ge0$ over the window days, the weighted mean and covariance of the asset returns are
\begin{equation}
\hat\mu(q)=\frac{\sum_s q_s r_s}{\sum_s q_s},\qquad
\hat\Sigma(q)=\frac{\sum_s q_s\,(r_s-\hat\mu(q))(r_s-\hat\mu(q))^\prime}{\sum_s q_s},
\end{equation}
with the maximum-likelihood normalisation. The pooled moments use $q_s=1$ and are written $\hat\mu,\hat\Sigma$. The regime moments use $q_s=\gamma_{s,k}$, with a floor of $10^{-12}$, and are written $\hat\mu_k,\hat\Sigma_k$. Write $\hat\Sigma=\mathrm{diag}(\hat\sigma)\hat R\,\mathrm{diag}(\hat\sigma)$ and $\hat\Sigma_k=\mathrm{diag}(\hat\sigma_k)\hat R_k\,\mathrm{diag}(\hat\sigma_k)$ for the decompositions into standard deviations and correlations of the raw returns. The three static models assign to regime component $k$ the covariance
\begin{align}
\text{A (pooled normal):}&\quad \Sigma^{A}_k=\hat\Sigma,\\
\text{B (regime volatility):}&\quad \Sigma^{B}_k=\mathrm{diag}(\hat\sigma_k)\,\hat R\,\mathrm{diag}(\hat\sigma_k),\\
\text{D (regime volatility and correlation):}&\quad \Sigma^{D}_k=\hat\Sigma_k.
\end{align}
Model A has the same covariance in both components, so it is a single normal distribution. Model B lets the volatilities differ by regime and holds the correlation matrix at its pooled value. Model D lets every volatility and every correlation differ by regime. The portfolio variance in component $k$ is $s_k^2=w^\prime\Sigma_k w$, and every component has the same mean $m=w^\prime\hat\mu$, the pooled window mean (the ``pooled'' mean treatment used in the ladder tables). The covariances are fixed within a 21-day block; only the mixture weights $p_{t,k}$ move from day to day. Two controls vary the mean: a regime-specific mean ($w^\prime\hat\mu_k$ in component $k$) and a zero mean. Appendix~\ref{app:mean} reports them, because they change the size of some contrasts.

\subsection{Models with GARCH margins: CCC, RSC and RSC\_P}\label{sec:garch}

These models replace the block-constant volatilities by GARCH(1,1) volatilities that update daily. For each asset $i$ let $x_{i,s}=100\,r_{i,s}$. On the window I estimate
\begin{equation}
x_{i,s}=\mu_i+\varepsilon_{i,s},\qquad \sigma_{i,s+1}^2=\omega_i+a_i\,\varepsilon_{i,s}^2+b_i\,\sigma_{i,s}^2,
\end{equation}
by Gaussian quasi-maximum likelihood, one asset at a time. Given the estimates, I re-filter the window with the recursion, starting from the sample variance of $\varepsilon_{i,\cdot}$ in the window, and form the standardised residuals $z_{i,s}=\varepsilon_{i,s}/\sigma_{i,s}$. The variance for the first day of the block comes from continuing the recursion, and for later days in the block it is updated each day from the realised return with the parameters held fixed. Write $\sigma_t=(\sigma_{1,t},\sigma_{2,t},\sigma_{3,t})^\prime$ for the volatility vector and $A_t=w\circ\sigma_t/100$, where $\circ$ is the element-wise product. The volatility vector is the same in every regime component: GARCH volatility in these models does not depend on the regime. The regime enters only through the correlation matrix of $z$ and the mixture weights.

The three models differ in the correlation matrix of the standardised residuals:
\begin{align}
\text{CCC:}&\quad R_1=R_2=\bar R,\ \text{a single component};\\
\text{RSC:}&\quad R_k=\bar R_k;\\
\text{RSC\_P:}&\quad R_k=\bar R\ \text{with the (SPY, IEF) entry replaced by } (\bar R_k)_{12}.
\end{align}
Here $\bar R$ is the correlation matrix of $z$ over the window with equal weights and $\bar R_k$ is the correlation matrix of $z$ with weights $\gamma_{s,k}$. In RSC\_P only the SPY--IEF correlation switches with the regime and the other two correlations stay pooled. I fixed this pair in advance from the stock--bond hypothesis, not by scanning. If a regime or hybrid matrix is not positive definite, it is repaired by clipping its eigenvalues and rescaling to a unit diagonal. The repair never fired, in the development data or in the hold-out run (counts in Chapter~\ref{ch:code}).

The forecast distribution of $y_t$ for these models is a mixture of normals with weights $p_{t,k}$ ($p_{t,k}=1$ for CCC), component means all equal to $m=w^\prime\hat\mu^{\mathrm{QML}}/100$ (the QML constants), and component variances $A_t^\prime R_k A_t$. The mean is the same in the two models of the primary test and is not regime-specific.

\subsection{Mixture value-at-risk and expected shortfall}\label{sec:mixvar}

For components with weights $p_k$, means $m_k$ and standard deviations $s_k$, the value-at-risk is the solution $q$ of
\begin{equation}
\sum_k p_k\,\Phi\!\left(\frac{q-m_k}{s_k}\right)=\alpha,
\end{equation}
found by root-finding, and the expected shortfall is
\begin{equation}
\ES=\frac1\alpha\sum_k p_k\Bigl[m_k\,\Phi(z_k)-s_k\,\phi(z_k)\Bigr],\qquad z_k=\frac{q-m_k}{s_k},
\end{equation}
where $\Phi$ and $\phi$ are the standard normal distribution and density functions. Both are returns, so both are negative at the levels used here. This is the Gaussian chain, which is the primary one. In the secondary Student-$t$ chain every component is a Student-$t$ distribution with common degrees of freedom $\nu$, scaled to the component's standard deviation. $\nu$ is estimated by maximum likelihood on the window's portfolio returns, standardised by an in-window model standard deviation; in the GARCH-margin models this is the constant-correlation standard deviation, so CCC, RSC and RSC\_P share one $\nu$. It is not comparable across the static and GARCH families.

\subsection{Benchmarks}\label{sec:bench}

Two other models serve as reference points. The first is a dynamic conditional correlation model (DCC; \citealp{engle2002}) built on the same GARCH(1,1) margins: the correlation matrix follows $Q_t=(1-\theta_1-\theta_2)\bar Q+\theta_1 z_{t-1}z_{t-1}^\prime+\theta_2 Q_{t-1}$, normalised to a correlation matrix, with $(\theta_1,\theta_2)$ estimated on the window by Gaussian quasi-maximum likelihood every 63 days and the recursion updated daily in between; if the estimates are not admissible, the model falls back to constant correlation (Chapter~\ref{ch:code}). The second is a GARCH(1,1) model with Student-$t$ errors fitted by maximum likelihood to the portfolio return itself (GARCH-t), which uses no information about the individual assets. Two practitioner baselines, both on the portfolio return, were added only after the hold-out was scored (Section~\ref{sec:holdout-posthoc}): historical simulation, the empirical 5\% quantile and the mean of the returns below it over the last 500 days; and a RiskMetrics exponentially weighted variance with decay 0.94, a zero mean and a normal distribution.

\subsection{What each comparison changes}\label{sec:ladder-design}

The thesis compares models in pairs that share the data, the regime signal and the estimation window. I do not call these pairs a nested sequence, because several pairs differ in more than one respect. Table~\ref{tab:differs} lists, for each comparison, the richer model and the baseline, what differs between them, and whether one model is a special case of the other.

\begin{table}[t]
\centering\setstretch{1}\scriptsize
\caption{\uc{What differs between the two models in each comparison}}\label{tab:differs}
\begin{tabular}{p{1.55cm}p{2.5cm}p{2.6cm}p{1.8cm}p{1.9cm}p{2.3cm}}
\toprule
Comparison (richer vs baseline) & Volatility & Correlation matrix & Mean & Components & Relation\\
\midrule
B vs A & regime-specific block SDs vs one pooled SD & pooled raw-return correlation in both & pooled window mean & 2 vs 1 & Nested in volatility\\
D vs B & same regime SDs & regime-specific vs pooled, raw returns & pooled window mean & same & Nested in the correlation matrix\\
CCC vs B & single-regime GARCH(1,1), daily update, vs regime-specific block SDs & standardised residuals vs raw returns & QML constant vs window mean & 1 vs 2 & Non-nested: changes four things\\
RSC vs D & single-regime GARCH(1,1) vs regime-specific block SDs & regime-specific, standardised residuals vs raw returns & QML constant vs window mean & 2 and 2 & Non-nested: model-class comparison\\
RSC vs CCC & identical GARCH & regime-specific vs pooled, same $z$ & identical & 2 vs 1 & Nested in the correlation matrix\\
RSC\_P vs CCC (primary) & identical GARCH & only the (SPY, IEF) entry regime-specific & identical & 2 vs 1 & Nested in one correlation entry\\
DCC vs CCC & same specification, fitted separately; own variance initialisation & recursion updated daily vs one pooled matrix per block & identical & 1 vs 1 & Parameter-nested ($\theta_1=\theta_2=0$), different code path\\
DCC vs RSC & as above & smooth recursion vs regime-weighted & identical & 1 vs 2 & Non-nested\\
\bottomrule
\end{tabular}
\tabnote{``Block SDs'' means standard deviations estimated on the window and held fixed for 21 days. B reduces to A when the two regime SDs are equal; D reduces to B and RSC to CCC when the regime correlation matrices equal the pooled one. Static arms use the pooled window sample mean and GARCH arms the QML constant mean. The DCC filter starts from the \texttt{arch} package's own variance initialisation and carries its correlation state forward daily, so it is not the CCC fit with a constant added.}
\end{table}

Three consequences follow. First, only the contrasts D versus B, RSC versus CCC and RSC\_P versus CCC change one thing, the correlation matrix, with the volatility model fixed. They are the only ones that support a statement about the correlation channel, and the statement holds only within the volatility model used. Second, the contrasts CCC versus B and RSC versus D change the volatility model, the correlation input, the mean estimator and the number of components together. They compare two classes of model, and they cannot be described as the effect of volatility dynamics: the GARCH models used here also remove the regime dependence from volatility. Third, because the margins of CCC, RSC and RSC\_P are the same, the primary test of the thesis, RSC\_P versus CCC, differs in two respects only: a mixture of two components against one, and a regime-specific SPY--IEF correlation entry against a pooled one.

Two families of comparisons were declared before they were scored. The first is the volatility-model contrast RSC versus D with a pooled mean, in the four ETF samples, with Holm's adjustment over the four (a gate: the ladder would become the main thesis if the contrast were negative in all four samples and significant after adjustment in at least two). The second is the confirmatory test RSC\_P versus CCC (Chapter~\ref{ch:eval}).
